% 389_Parabelflug_Darstellungen_En_ch.tex
% Johann Pascher, 6 Oct. 2026
% Parabolic flight, inertia and the representations of gravitation

\providecommand{\BM}{\textbf{[B]}}
\providecommand{\KM}{\textbf{[K]}}
\providecommand{\SM}{\textbf{[S]}}
\providecommand{\XM}{\textbf{[X]}}
\providecommand{\QM}{\textbf{[Q]}}

\chapter*{Parabolic Flight, Inertia and the Representations of Gravitation}
\addcontentsline{toc}{chapter}{Parabolic Flight, Inertia and the Representations of Gravitation}

\begin{abstract}
In parabolic flight, gravity disappears for the passengers. This document starts from that
question and reads it with the time-mass duality $\tilde T\cdot m=1$: free fall is the path of
stationary phase of a mode with position-dependent mass, and weight on the ground is the deviation
of the rest world line from that path. What cannot be transformed away are the tides. Gravitation
appears in three representations: as space curvature, as a change of mass or clock rate, and as
fractal path lengthening. They are mathematical models that must be equivalent by virtue of the
duality, not statements about whether space is ``really'' flat or curved. For light, clock rate
and path are complementary halves: each alone gives $0.875''$ at the solar limb, both together
$1.75''$ and $\gamma=1$. The purely conformally flat metric gives no deflection. With the
logarithmic form of the mass change, $\beta=1$ follows as well, and with it the perihelion shift
of GR. At the horizon, $T_0$ limits the comparison with a distant reference, not the free
crossing. A separate quantisation of gravitation is not needed.
\end{abstract}

% ============================================================
\section*{Status markers}
% ============================================================
\begin{center}
\begin{tabular}{@{}lp{11cm}@{}}
\textbf{[B]} & Proved algebraically. \\[2pt]
\textbf{[K]} & Confirmed numerically (verification script, comparison with measured values). \\[2pt]
\textbf{[S]} & Setting, open or speculative. \\[2pt]
\textbf{[X]} & Refuted or not viable. \\[2pt]
\textbf{[Q]} & Convention. \\
\end{tabular}
\end{center}
Measured values: Pound--Rebka 1960, Cassini (Bertotti et al.\ 2003), IAU and CODATA 2022
constants. Verification script: \texttt{pruef\_389\_parabelflug.py}.

% ============================================================
\section{The question}
% ============================================================

In parabolic flight the aircraft flies a ballistic parabola for about twenty seconds. The
passengers float, although the Earth attracts them as before. On the ground, by contrast, they
feel their weight, although they are not moving. The question is what this weight is, if it
disappears in free fall.

The answer this document works out: the force we feel as weight exists only because we deviate
from the path we would take without support. The ground forces us onto a path that is accelerated
exactly opposite to the path of fall. In parabolic flight this deviation drops out, and with it
the weight. The following sections recalculate this in the representation of FFGFT and place it
among the different readings of gravitation.

% ============================================================
\section{Free fall in the mass representation}
% ============================================================

In the representation ``space flat, mass variable'' (Doc.~078) a mode in the potential $\Phi$
carries the position-dependent mass $m(x)$. Its action is that of a clock ticking at the rate
$m(x)c^2$:
\begin{equation}
L=-m(x)\,c^2\sqrt{1-v^2/c^2},\qquad m(x)=m_0\,e^{\Phi(x)/c^2}.
\end{equation}
Here $m(x)$ is the mass as compared from the distant reference. A clock at $\Phi$ ticks in that
reference's time at
\begin{equation}
\omega(\Phi)=\omega_0\sqrt{g_{00}}\approx\omega_0\Bigl(1+\frac{\Phi}{c^2}\Bigr),\qquad
m(x)=m_0\sqrt{g_{00}},\qquad g_{00}=\Bigl(\frac{m}{m_0}\Bigr)^2=e^{2\Phi/c^2}\quad\BM,
\end{equation}
since with $\tilde T\cdot m=1$ the mass is the frequency of the mode. $m(x)$ is the Killing energy
of the mode; measured with a co-moving clock the mass remains $m_0$.
For small velocity the Euler--Lagrange equation gives
\begin{equation}
\ddot{\vec x}=-c^2\,\nabla\ln m=c^2\,\nabla\ln\lambda_4=-\nabla\Phi\quad\BM,
\end{equation}
with the mode length $\lambda_4=2\pi/m$ (Doc.~314). Without a metric and without Christoffel
symbols, the fall follows from the gradient of the mass alone: bodies fall towards where their
mass is smaller and their mode length larger. Since $\nabla\ln m$ is the same for every mode, all
bodies fall alike; in this form the equivalence principle is built in.

At the Earth's surface $g/c^2=1.09\times10^{-16}$ per metre \KM. This is the same number by which a
clock runs faster per metre of height. In this representation, gravitational acceleration and
clock rate are a single quantity.

\paragraph{Sign.} The direction of the mass change is fixed by the redshift. Light rising in a
gravitational field arrives redshifted; the Pound--Rebka experiment measured this over $22.5$\,m
as $(2.57\pm0.26)\times10^{-15}$, with $2.46\times10^{-15}$ expected \KM. In the mass
representation this means: the mode below, compared from the upper reference, has the smaller
mass; $m$ grows with $\Phi$. With this sign the body falls downwards. The sign is therefore not an
open choice but determined by measurement.

% ============================================================
\section{Rest as a counter-path}
% ============================================================

The phase of a mode along a path is $\varphi=-\int m(x)\,c^2\,d\tau$. Free fall is the path on
which this phase is stationary, $\delta\varphi=0$. The rest world line $\vec x=\text{const}$ is
not: its variation does not vanish but leaves the remainder
\begin{equation}
\frac{\partial L}{\partial \vec x}\Big|_{v=0}=-m\,\nabla\Phi=-m\,\vec g\quad\BM.
\end{equation}
For a body to be at rest nonetheless, the ground must balance exactly this remainder with
$+m\vec g$ upwards. That is the weight. It is not an additional force on top of gravity, but the
measure of how far the rest world line deviates from the path of stationary phase. In the language
of the mode it is a phase detuning: the mode at rest runs against the rate it would have in free
fall.

In parabolic flight the aircraft follows the path of fall. The deviation vanishes, and with it
the weight. What the passengers experience as heaviness on the ground and as weightlessness in the
aircraft is the same path, once with and once without the counter-path. As a formula:
\begin{equation}
\vec a_{\text{fall}}=c^2\,\nabla\ln\lambda_4,\qquad
\vec a_{\text{inert}}=-c^2\,\nabla\ln\lambda_4=-\vec a_{\text{fall}}\quad\BM.
\end{equation}
On the ground the two cancel; in parabolic flight the counter-path is zero. In this sense the
inertial force exists only because we have an exactly opposite path.

% ============================================================
\section{What cannot be transformed away}
% ============================================================

The gradient $\nabla\ln m$ can be transformed away locally: in the freely falling frame it is
zero, as are the Christoffel symbols in the local inertial frame of GR. The second derivative
cannot be transformed away. The tides are the Hessian
\begin{equation}
\partial_i\partial_j\Phi=c^2\,\partial_i\partial_j\ln m,\qquad
\text{eigenvalues }(-2,\,1,\,1)\,\frac{GM}{r^3}\ \text{on the axis},\qquad
\nabla^2\ln m=\frac{4\pi G\rho}{c^2}\quad\BM.
\end{equation}
In vacuum the trace is zero; inside a mass it gives the Poisson equation, now as an equation for
the mass field. Outside, $m=m_0\,e^{-GM/(rc^2)}\approx m_0\bigl(1-GM/(rc^2)\bigr)$; the gradient
$\partial_r\ln m=GM/(r^2c^2)$ can be transformed away, the second derivative
$c^2\,\partial_r^2\ln m=-2GM/r^3$ cannot \BM. At the Earth's surface $2GM/r^3=3.08\times10^{-6}\,\text{s}^{-2}$ \KM. This
remainder persists in parabolic flight as well: two bodies at different places in the aircraft do
not fall exactly alike. What the mass representation calls ``curvature'' is the curvature of the
mass field, not that of a spatial geometry.

% ============================================================
\section{Three representations}
% ============================================================

A position-dependent mass is equivalent to a position-dependent time, since
$\tilde T\cdot m=1$. What one calls it is a choice of representation:
\begin{itemize}[leftmargin=*,itemsep=1pt]
\item \textbf{Space curvature:} mass constant, space and time curved, as in GR.
\item \textbf{Change of mass or clock rate:} space flat, mass and clock rate position-dependent
(Doc.~078).
\item \textbf{Fractal path lengthening:} space flat, propagation paths geometrically stretched
(Docs.~026, 041, 182).
\end{itemize}
These are mathematical models that must be equivalent by virtue of the duality. Ontologically, one
can say neither that space is flat nor that it is curved; mixed forms are conceivable. All that
matters is that the observables agree.

\paragraph{Light as the test.} For clocks, redshift and slow bodies the clock rate suffices. Light
tests more, because it also feels the spatial geometry. In Fermat form light propagates through a
refractive index $n=1+(1+\gamma)\,GM/(rc^2)$; the deflection at the solar limb and the Shapiro
delay scale with $(1+\gamma)/2$:
\begin{center}
\begin{tabular}{lcc}
\toprule
Representation & $\gamma$ & Deflection at the solar limb \\
\midrule
purely conformally flat metric (Doc.~143) & $-1$ & $0$ \\
clock rate only (Einstein 1911) & $0$ & $0.875''$ \\
clock rate and fractal path lengthening (Doc.~308) & $1$ & $1.751''$ \\
\bottomrule
\end{tabular}
\end{center}
\KM. Cassini measures $\gamma-1=(2.1\pm2.3)\times10^{-5}$; only the last row is compatible with the
measurement. Doc.~308 computes the same by ray tracing: clock rate and path lengthening give
$0.875''$ each, together $1.750''$. They are complementary halves, not alternatives.

This sharpens the three representations. ``Space flat, mass variable'' on its own carries only the
temporal half. The flat representation becomes complete only with the fractal path lengthening,
which carries the spatial half. The space curvature of GR is then equivalent to clock rate and
path together. The temporal half is measured by clocks and is the same in every representation;
the spatial half can be written as curvature, as path lengthening, or as a mixture of both. The
freedom of representation therefore lies in the spatial half, and the condition for every mixed
form is $\gamma=1$.

The conformally flat metric $g_{\mu\nu}=\eta_{\mu\nu}(1-2\xi E_{\text{field}}/E_P)$ of Doc.~143 is
none of these representations: it stretches space and time by the same factor, null geodesics
remain unchanged, and light propagates as in flat space. This is the Nordström case (Doc.~143). The same holds for the conformal coupling
$g_{\mu\nu}\to\Omega^2g_{\mu\nu}$ of the matter fields to the time field (A142): the Maxwell field is
conformally invariant in four dimensions, and the coupling leaves light unaffected \BM. For massive
matter it supplies the clock rate, i.e.\ the temporal half; the spatial half has to come from the
path lengthening (A142).

% ============================================================
\section{Second order: the perihelion shift}
% ============================================================

Light deflection tests the first order in the potential. The perihelion shift additionally tests
the second order of the time component, the parameter $\beta$ in $g_{00}=1-2U+2\beta U^2$ with
$U=GM/(rc^2)$. The factor relative to GR is $(2+2\gamma-\beta)/3$. In the mass representation
$g_{00}=(m/m_0)^2$, and $\beta$ depends on the form of $m(\Phi)$:
\begin{center}
\begin{tabular}{lccc}
\toprule
Form & $\beta$ & with $\gamma$ & Mercury ($''$/cent.) \\
\midrule
$m=m_0\,e^{\Phi/c^2}$ & $1$ & $1$ & $42.98$ (GR) \\
$m=m_0\,(1+\Phi/c^2)$ & $\tfrac12$ & $1$ & $50.1$ \\
Nordström & $\tfrac12$ & $-1$ & $-7.2$ \\
clock rate only & $1$ & $0$ & $14.3$ \\
\bottomrule
\end{tabular}
\end{center}
\KM. The logarithmic form $\ln(m/m_0)=\Phi/c^2$ is the same one that in Section~2 gives the fall
$\ddot{\vec x}=-\nabla\Phi$ exactly and not merely approximately. With it, and with the spatial
half from Doc.~308, the flat representation gives the perihelion shift of GR. The logarithmic form
itself is set, not derived from $\tilde T\cdot m=1$ \SM; the linear form is excluded by the
Mercury measurement.

% ============================================================
\section{The relation \texorpdfstring{$\lambda_4\cdot m=2\pi$}{lambda4 m = 2 pi} holds locally}
% ============================================================

The position relation $\lambda_4\cdot m=2\pi$ (Doc.~314) connects proper quantities: the mode
length as a co-moving rod measures it, and the mass as a co-moving clock measures it. In
coordinates of the distant reference it does not hold unchanged once the spatial half is written
as curvature, because coordinate lengths then carry a factor of the spatial metric. In the
isotropic form of GR,
\begin{equation}
g_{ij}=\Bigl(1-\frac{2\Phi}{c^2}\Bigr)\delta_{ij},\qquad
\ell_{\text{coord}}=\ell\,\Bigl(1+\frac{\Phi}{c^2}\Bigr),\qquad
\lambda_{4,\text{coord}}=\frac{2\pi}{m(x)}=\lambda_4\Bigl(1-\frac{\Phi}{c^2}\Bigr)\quad\BM.
\end{equation}
Seen from the distant reference, spatial lengths deep in the potential appear shortened and the
temporal mode length lengthened: the two factors run against each other. In the purely
conformally flat metric both would scale alike with $1/m$, and light would remain undeflected.
The spatial factor is exactly the half that light additionally feels. In the representation with
fractal path lengthening it sits in the stretched path, as a refractive contribution
$n-1=-\Phi/c^2$ added to the clock rate. The
gradient $\nabla\ln\lambda_4$ in Section~2 is accordingly a local quantity: in the freely falling
frame it vanishes, in the frame at rest it is $g/c^2$.

% ============================================================
\section{Horizon: comparison and crossing}
% ============================================================

Doc.~306 describes the horizon as an extremum: there $T$ reaches the value $T_0=\xi\,t_P$ and the
energy per excitation the value $E_P/\xi$. It adds that the ``slowest time'' at the horizon is a
statement about the comparison, not about experience. The calculation confirms this.

A static observer at $r$ measures, for a mode of energy $E$ (referred to the distant reference),
the local energy $E/\sqrt{1-r_s/r}$; it diverges at the horizon \BM. The bound $E\le E_P/\xi$
cuts this approach off at the proper distance $\rho_{\min}=2\,r_s\,\xi\,E/E_P$ \BM. A freely
falling observer, by contrast, measures $E/(1+\sqrt{r_s/r})$ for incoming light, hence $E/2$ at the
horizon, which is finite \BM. For this observer there is nothing to bound at the horizon.

The bound $T_0$ therefore concerns the static reading, the comparison with a distant reference.
The free crossing is not affected; this is the same distinction as between rest and fall in
parabolic flight. Whether $T_0$ additionally bounds the radius in the interior from below requires
a dynamics of the interior that the corpus does not supply. The physical validity in the
black-hole regime remains open \SM\ (Doc.~384).

% ============================================================
\section{Two masses, two readings of \texorpdfstring{$G$}{G}}
% ============================================================

Two masses at the upper end are easily confused:
\begin{itemize}[leftmargin=*,itemsep=1pt]
\item $M_{\text{coll}}=m_P/\sqrt2$: here the Schwarzschild radius equals the reduced Compton length
\BM\ (Docs.~325, 329).
\item $m_P/\xi=7500\,m_P\approx1.63\times10^{-4}$\,kg: the dual upper bound per excitation,
belonging to $L_0=\xi\,\ell_P$ \KM\ (Docs.~290, 306).
\end{itemize}
Both are correct but different; their ratio is $7500\sqrt2\approx10\,607$. According to Doc.~306,
both bounds hold per excitation, not per aggregate: a star is not a single winding, and
$\hbar/(Mc^2)$ for a stellar mass is not a state of the relation.

The same applies to $G$. With the measured, fixed $G$, $m\cdot G=\xi^2/4$ holds for the
characteristic mass $m_{\text{char}}$ of the anchor (R58, A142). For arbitrary $m$ the relation
holds only if $G$ runs along, as Doc.~350 reads it for an evaporating hole. The statement ``no mass
singularity with finite $G$'' is, in the first reading, a consequence of the definition; in the
second, a statement about the dynamics.

% ============================================================
\section{No separate quantisation of gravitation}
% ============================================================

In FFGFT gravitation is not an additional field and needs no graviton. It is already contained in
the time-field geometry and shows itself as an effect in the redshift, read as space curvature, as
mass change or as fractal path lengthening. A separate quantisation of gravitation is therefore
not required (R58, A142). What remains open is the transfer of the classical equations into loop
corrections for situations beyond the Schwarzschild limit (A142).

% ============================================================
\section{What remains open}
% ============================================================

\begin{itemize}[leftmargin=*,itemsep=1pt]
\item the logarithmic form $\ln(m/m_0)=\Phi/c^2$, which gives $\beta=1$, as a consequence of the
duality rather than only as a setting \SM;
\item the origin of the fractal path lengthening as the spatial half from the structure of the
carrier; Doc.~308 gives the K3 area invariance as the reason, confirmed numerically but not proved
\KM;
\item the dynamics inside a horizon and the validity of $m\cdot G=\xi^2/4$ beyond the anchor \SM.
\end{itemize}

% ============================================================
\section{Conclusion}
% ============================================================

In parabolic flight it is not gravitation that disappears but the counter-path. Free fall is the
path of stationary phase of a mode with position-dependent mass, weight on the ground is the
deviation of the rest world line from that path, and only the tides persist in every frame. Space
curvature, mass change and fractal path lengthening are equivalent models of the same situation,
not statements about the nature of space. For light, clock rate and path are complementary halves;
only both together give $1.75''$ and $\gamma=1$. With the logarithmic form of the mass change the
perihelion shift of GR follows as well. At the horizon, the same distinction between rest and fall
separates the comparison with a distant reference from the free crossing.

\vspace{1em}
\noindent\textit{Verification script:} \texttt{2/python/Dok389\_Skripte/pruef\_389\_parabelflug.py}
--- 34/34.
